\begin{figure}[htbp]
\centering
\begin{tikzpicture}[
    font=\small,
    caja/.style={draw=#1, fill=#1!8, rounded corners=3pt, align=center,
                 minimum width=3.4cm, minimum height=1.15cm, inner sep=3pt},
    caja/.default=subsectionblue,
    hist/.style={draw=gray!60, fill=gray!8, rounded corners=2pt, align=left,
                 font=\scriptsize\ttfamily, inner sep=4pt, minimum width=4.6cm},
    flecha/.style={-{Stealth[length=2mm]}, thick, draw=gray!70!black},
    et/.style={font=\scriptsize, text=gray!55!black, align=center, inner sep=2pt},
]

% --- Actores ---------------------------------------------------------------
\node[caja]        (ag)  at (0,0)    {\texttt{agent.py}\\[-1pt]\scriptsize bucle agéntico};
\node[caja=ugrred] (llm) at (7.6,0)  {Modelo\\[-1pt]\scriptsize \texttt{qwen3:8b}};
\node[caja=gray]   (fn)  at (0,-3.8) {Función Python\\[-1pt]\scriptsize \texttt{ejecutar\_tool}};

% --- Ciclo -----------------------------------------------------------------
\draw[flecha] ([yshift=3mm]ag.east) -- node[et, above]
      {1. mensajes + esquemas\\de las 6 herramientas} ([yshift=3mm]llm.west);
\draw[flecha] ([yshift=-3mm]llm.west) -- node[et, below]
      {2. \texttt{tool\_calls}: nombre\\y argumentos en JSON} ([yshift=-3mm]ag.east);
\draw[flecha] (ag) -- node[et, left] {3. ejecuta} (fn);
\draw[flecha] ([xshift=6mm]fn.north) -- node[et, right]
      {4. resultado como\\mensaje de rol \texttt{tool}} ([xshift=6mm]ag.south);

\draw[flecha, draw=subsectionblue] (llm.north) .. controls (7.6,2.0) and (0,2.0) ..
      node[et, above, text=subsectionblue, pos=0.5] {se repite hasta que el modelo responde sin pedir herramientas (máx.\ 8 vueltas)}
      (ag.north);

% --- Historial resultante --------------------------------------------------
\node[hist, anchor=north west] (h) at (4.0,-2.6)
      {\textcolor{gray}{system}: reglas y esquema\\
       \textcolor{gray}{user}: ¿cuánto gasté en gasolina?\\
       \textcolor{gray}{assistant}: \textit{tool\_calls}\\
       \textcolor{gray}{tool}: \{"filas": [[84.2]]\}\\
       \textcolor{gray}{assistant}: Has gastado 84,20 €};
\node[et, anchor=south west] at (4.0,-2.5) {El historial que se reenvía en cada vuelta};

\end{tikzpicture}
\caption{Ciclo de \textit{tool calling}. El modelo nunca ejecuta nada: devuelve
el nombre de la herramienta y sus argumentos, el \emph{backend} ejecuta la
función en Python y le devuelve el resultado como un mensaje más del historial.
Con ese dato ya concreto, el modelo redacta la respuesta final.}
\label{fig:flujo_agente}
\end{figure}
